\chapter{Systematic Credit and Bond-ETF Arbitrage}\label{ch:s2:systematic-credit-and-bond-etf-arbitrage}

A bond ETF trades every second, while most of its bonds do not trade for days. When the two disagree, authorised participants create or redeem shares,
and the discount says how illiquid the bonds really are. In March 2020, Haddad, Moreira and Muir found, liquid bond ETFs traded at large discounts to
their net asset values as investors sold what they could to raise cash. The same market has systematic factors. Houweling and van Zundert found size,
low-risk, value and momentum premia in corporate bonds, with low correlations between them. On this chapter's synthetic bond market, low risk earns a
Sharpe ratio of 0.95, value 0.46 and momentum 0.14. In a planted stress, the ETF's price falls 5.0\% below its stale NAV, 3.0\% below the bonds' true
value, and an authorised participant can profit on 28 days. The build is \texttt{firm.syscredit}.

\section{Factors in credit}

\begin{definition}[Credit factor]\label{def:s2:systematic-credit-and-bond-etf-arbitrage:factor}
A \emph{credit factor}\index{credit factor} is a characteristic of corporate bonds (spread against peers of similar rating and duration, recent spread
change, duration or rating, issuer size) that sorts them into portfolios whose returns in excess of Treasuries differ persistently, after adjusting for
their spread duration.
\end{definition}

Houweling and van Zundert found that size, low-risk, value and momentum portfolios generated economically meaningful and statistically significant
alphas in the corporate bond market. They held up after transaction costs and within liquid bonds, and combined well because their correlations were
low. Measuring factors in bonds has two difficulties that equities do not share. A bond's return is mostly its spread duration times its spread
change, so every signal must be compared across similar durations. And bond prices are stale, so a signal read from last week's quotes may already be
out of date.

\texttt{firm.syscredit} builds 300 bonds with durations from one to twelve years over eight years
(\cref{lst:s2:systematic-credit-and-bond-etf-arbitrage:factors}). Each spread moves with a market factor, an issuer trend that persists for months,
a gap from fair value that reverts over a year, and noise. Expected excess returns rise with the square root of duration, so short, safe bonds earn
more per unit of risk. Each factor book buys the top fifth and sells the bottom fifth of bonds by its signal, weighted so that each side has a similar
spread duration. Value is the spread above the fitted spread-duration line less its three-month change, which keeps value from being a list of recent
losers. Momentum is the six-month spread tightening. Low risk buys short bonds and sells long ones at equal risk.

\begin{center}
\small
\begin{tabular}{@{}lcc@{}}
\toprule
eight years, monthly rebalancing & Sharpe ratio & correlation with value\\
\midrule
value & 0.46 & 1\\
momentum & 0.14 & $-0.09$\\
low risk & 0.95 & $-0.03$\\
equal-risk combination & 0.94 &\\
\bottomrule
\end{tabular}
\end{center}

Low risk dominates the combination (\cref{fig:s2:systematic-credit-and-bond-etf-arbitrage:factors}). Momentum is weak because the value gap, which
reverts, drives part of each bond's six-month move. The two signals fight over the same spread changes. Houweling and van Zundert found both
positive with real data. A model in which one reverts and the other trends at overlapping horizons gives neither much room.

\begin{omfigure}
\begin{tikzpicture}
\begin{axis}[width=11cm,height=5.4cm,xlabel={\small year},ylabel={\small cumulative (\%, at 10\% vol)},tick label style={font=\footnotesize},
  xmin=0,xmax=7.5,ymin=-20,ymax=85,grid=major,grid style={gray!25},table/col sep=comma,
  legend style={font=\scriptsize,at={(0.5,-0.3)},anchor=north},legend columns=4]
\addplot[omDef,thick] table[x=year,y=value]{figdata/strategies-2/19-systematic-credit-and-bond-etf-arbitrage/factors.csv};
\addplot[gray,thick] table[x=year,y=momentum]{figdata/strategies-2/19-systematic-credit-and-bond-etf-arbitrage/factors.csv};
\addplot[omThm,thick] table[x=year,y=lowrisk]{figdata/strategies-2/19-systematic-credit-and-bond-etf-arbitrage/factors.csv};
\addplot[black,very thick] table[x=year,y=combined]{figdata/strategies-2/19-systematic-credit-and-bond-etf-arbitrage/factors.csv};
\legend{value,momentum,low risk,combined}
\end{axis}
\end{tikzpicture}

\omcaption{Three synthetic \omterm{def:s2:systematic-credit-and-bond-etf-arbitrage:factor}{credit factor} books and their equal-risk combination, each scaled to 10\% volatility, cumulative sums of daily returns.
Data: \texttt{s2\_syscredit.market}.}\label{fig:s2:systematic-credit-and-bond-etf-arbitrage:factors}
\end{omfigure}

\section{Bond ETFs and create-redeem}

\begin{definition}[ETF create-redeem arbitrage]\label{def:s2:systematic-credit-and-bond-etf-arbitrage:createredeem}
\emph{ETF create-redeem arbitrage}\index{ETF create-redeem arbitrage} is the trade by which an authorised participant (Book~1, chapter~14) delivers a
creation basket to receive new ETF shares when the shares trade above the basket's value, or delivers shares to receive the basket when they trade
below, keeping the gap less its costs.
\end{definition}

For a bond ETF the basket is the hard part. Koont, Ma, Pastor and Zeng found that corporate bond ETFs choose creation and redemption baskets that hold
cash and only a subset of the index bonds, the more so when the bonds are illiquid. They adjust the baskets to correct imbalances while still
allowing arbitrage.
Basket inclusion made bonds more liquid in general but less liquid when creations and redemptions were most unbalanced, as in the COVID-19 crisis.

\begin{definition}[ETF discount]\label{def:s2:systematic-credit-and-bond-etf-arbitrage:discount}
An \emph{ETF discount}\index{ETF discount} is the amount by which an ETF's market price is below its published net asset value; for a bond ETF part of
it is stale NAV, computed from bond prices that have not traded, and part is selling pressure on the shares that the creation-redemption mechanism has
not absorbed.
\end{definition}

\section{Portfolio trades}

A portfolio trade (Book~2, chapter~22) buys or sells a whole list of bonds at once, often an ETF's basket, at a single negotiated price. Dealers price
it against the ETF and the bonds' quotes, and can hedge it with the ETF itself. When the ETF trades at a discount, a client selling a portfolio of
bonds pays for the dealer's cost of redeeming or holding them. Providing liquidity to portfolio trades is a way of being the authorised
participant's counterparty without the creation unit.

\section{Discounts in stress}

The synthetic ETF holds all 300 bonds. Its NAV uses the bonds' last traded prices; each bond trades on a quarter of days. Its market price follows the
bonds' true value. Six years in, a planted stress widens market spreads by 150 basis points over three weeks, and the basket loses 9.8\%. For the same
weeks the ETF price falls a further 3\% below true value, as investors sell ETF shares for cash, and recovers over two months
(\cref{fig:s2:systematic-credit-and-bond-etf-arbitrage:etf}).

In normal times the discount averages 2.7 basis points with a standard deviation of 21.6. At the stress's worst, fourteen trading days after it began, the ETF traded
504 basis points below NAV. Of that, the stale NAV stood 215 basis points above the bonds' true value, and the price stood 300 basis points below it.
Only the second part is an arbitrage. An authorised participant who buys shares, redeems them and sells the bonds at their true value pays a cost of
1.5\% in the stress (0.3\% normally). The trade was profitable on 28 days, by up to 1.63\% of the value redeemed, and on no day outside the stress.

\begin{omfigure}
\begin{tikzpicture}
\begin{axis}[width=11cm,height=5.4cm,xlabel={\small trading days from the stress's start},ylabel={\small value (100 the day before)},
  tick label style={font=\footnotesize},xmin=-10,xmax=64,ymin=86,ymax=102,grid=major,grid style={gray!25},table/col sep=comma,
  legend style={font=\scriptsize,at={(0.97,0.97)},anchor=north east}]
\addplot[black,thick] table[x=day,y=true]{figdata/strategies-2/19-systematic-credit-and-bond-etf-arbitrage/etf.csv};
\addplot[omDef,thick,dashed] table[x=day,y=nav]{figdata/strategies-2/19-systematic-credit-and-bond-etf-arbitrage/etf.csv};
\addplot[omThm,thick] table[x=day,y=price]{figdata/strategies-2/19-systematic-credit-and-bond-etf-arbitrage/etf.csv};
\legend{bonds' true value,NAV from stale prices,ETF price}
\end{axis}
\end{tikzpicture}

\omcaption{The synthetic bond ETF through a three-week stress: the true value of its 300 bonds, the NAV computed from their last traded prices, and the
ETF's market price, which falls below both. Data: \texttt{s2\_syscredit.market}.}\label{fig:s2:systematic-credit-and-bond-etf-arbitrage:etf}
\end{omfigure}

Haddad, Moreira and Muir read the March 2020 discounts the same way. Investors tried to sell their safer, more liquid holdings to raise cash, and the
discounts were larger for Treasury, municipal and investment-grade ETFs than for high-yield ones. The disruptions reversed almost as fast as they
appeared, after the Federal Reserve announced purchases of corporate bonds on 23 March and 9 April 2020. Which price was right? The NAV was stale, the
ETF price was pushed, and the bonds' true value was the one nobody observed.

\section{Strategy files}

\begin{strategyfile}{Credit value and momentum}\label{strat:s2:systematic-credit-and-bond-etf-arbitrage:factors}
\sfield{Who pays you, and why} Investors constrained by ratings and benchmarks, who leave cheap bonds cheap and chase recent winners slowly.
\sfield{Instruments and venues} Corporate bonds; CDS as a liquid substitute.
\sfield{Signal} Spread against peers of similar rating and duration; spread momentum; hedged to duration.
\sfield{Sizing and execution} Long-short or tilts against a benchmark; monthly.
\sfield{Costs} Bond bid-ask, large for small and old issues.
\sfield{How it dies} Costs; value that is default risk; momentum crashes at turning points.
\sfield{Horizon, capacity, infrastructure} Months; a bond database with stale-price flags.
\sfield{Backtest honestly} Traded prices, not matrix prices; costs by liquidity bucket.
\sfield{Sources} Houweling and van Zundert (2017); this chapter: value 0.46, momentum 0.14.
\end{strategyfile}

\begin{strategyfile}{Low-risk credit}\label{strat:s2:systematic-credit-and-bond-etf-arbitrage:lowrisk}
\sfield{Who pays you, and why} Investors who reach for yield in long, risky bonds and leave short, safe ones cheap per unit of risk.
\sfield{Instruments and venues} Short-dated, higher-rated bonds, levered or against long, risky ones.
\sfield{Signal} Duration and rating.
\sfield{Sizing and execution} Equal risk on each side.
\sfield{Costs} Funding for leverage.
\sfield{How it dies} A rally in risky credit; funding costs.
\sfield{Horizon, capacity, infrastructure} Years.
\sfield{Backtest honestly} Leverage at realistic funding rates.
\sfield{Sources} Houweling and van Zundert (2017); this chapter: 0.95.
\end{strategyfile}

\begin{strategyfile}{ETF premium-discount arbitrage}\label{strat:s2:systematic-credit-and-bond-etf-arbitrage:etf}
\sfield{Who pays you, and why} ETF holders who sell shares for cash faster than the bonds can be sold.
\sfield{Instruments and venues} ETF shares; creation and redemption with the sponsor as an authorised participant; the bonds.
\sfield{Signal} The price against the bonds' true value, not against a stale NAV.
\sfield{Sizing and execution} Redeem when the gap exceeds the cost of selling the basket; baskets as the sponsor sets them.
\sfield{Costs} Selling illiquid bonds in stress.
\sfield{How it dies} A discount that is all stale NAV; baskets changed to hold bonds that cannot be sold.
\sfield{Horizon, capacity, infrastructure} Days; authorised-participant status and bond trading desks.
\sfield{Backtest honestly} Estimate true value from traded prices; include basket rules.
\sfield{Sources} Haddad, Moreira and Muir (2020); Koont, Ma, Pastor and Zeng (2022); this chapter: 28 profitable days in the stress.
\end{strategyfile}

\begin{strategyfile}{Portfolio-trade liquidity provision}\label{strat:s2:systematic-credit-and-bond-etf-arbitrage:portfolio}
\sfield{Who pays you, and why} Clients who pay to trade a whole list of bonds at once.
\sfield{Instruments and venues} Portfolio trades priced against the ETF; the ETF as a hedge.
\sfield{Signal} The list's price against the ETF and the bonds' quotes.
\sfield{Sizing and execution} Hedge with the ETF at once, work out of the bonds over days.
\sfield{Costs} The bonds' bid-ask as they are worked out.
\sfield{How it dies} Lists of bonds that only the dealer holds; ETF and bonds diverging.
\sfield{Horizon, capacity, infrastructure} Days; balance sheet.
\sfield{Backtest honestly} Execution of each bond at the prices that were available.
\sfield{Sources} No performance figure verified.
\end{strategyfile}

\section{Tutorial: which price is right}\label{tut:s2:systematic-credit-and-bond-etf-arbitrage:tut}

\begin{tutorial}
\textbf{Goal.} Build a bond universe with factor structure and stale prices, trade \omterm{def:s2:systematic-credit-and-bond-etf-arbitrage:factor}{credit factors}, and run the ETF and its authorised participant
through a stress. \textbf{End state:} the table and the two figures.
\begin{tutsteps}
\item \textbf{Factor books}.
\omcode{code/firm/syscredit/firm_syscredit.py}{91}{119}{Value, momentum and low-risk books, duration-balanced, monthly.}\label{lst:s2:systematic-credit-and-bond-etf-arbitrage:factors}
\item \textbf{The authorised participant}.
\omcode{code/firm/syscredit/firm_syscredit.py}{122}{132}{Redeeming at a discount beyond the cost of selling the bonds.}\label{lst:s2:systematic-credit-and-bond-etf-arbitrage:ap}
\item \textbf{Run} \texttt{factors()}, \texttt{etf()} and \texttt{fig\_syscredit.py}.
\end{tutsteps}
\textbf{What to change next.} Estimate the bonds' true value from the ETF and recent trades and trade the difference; let the basket omit the least
liquid bonds; add a Fed-purchase announcement that closes the dislocation in a day.
\end{tutorial}

\section{Build: systematic credit}\label{bld:s2:systematic-credit-and-bond-etf-arbitrage:syscredit}

\begin{build}
\bfield{Purpose} A synthetic bond panel with \omterm{def:s2:systematic-credit-and-bond-etf-arbitrage:factor}{credit factors} and stale prices, a bond ETF with NAV and price, and create-redeem arbitrage.
\bfield{Interface} \texttt{SysCreditConfig(\dots)}, \texttt{simulate\_bonds(cfg)}, \texttt{factor\_book(sim, cfg, name)},
\texttt{ap\_arbitrage(sim, cfg)}.
\bfield{Rules} Stale prices update only on trading days; NAV from stale prices; arbitrage against true value less cost.
\bfield{Acceptance tests} \texttt{code/firm/syscredit/tests/}: stale prices move only on trades and NAV is their mean; no arbitrage without a
dislocation; factor books are long-short and finite.
\bfield{Stretch} Rating migration; defaults; basket choice.
\end{build}

\begin{omsources}
\item P.~Houweling and J.~van Zundert, ``Factor investing in the corporate bond market'', \emph{Financial Analysts Journal} 73(2), 2017.
\item V.~Haddad, A.~Moreira and T.~Muir, ``When selling becomes viral: disruptions in debt markets in the COVID-19 crisis and the Fed's response'',
NBER Working Paper 27168, 2020.
\item N.~Koont, Y.~Ma, L.~Pastor and Y.~Zeng, ``Steering a ship in illiquid waters: active management of passive funds'', NBER Working Paper 30039,
2022.
\end{omsources}

\section{Exercises}

\begin{exercise}[$\star$]\label{exo:s2:systematic-credit-and-bond-etf-arbitrage:1}
An ETF trades 504 basis points below NAV; the NAV is 215 basis points above the bonds' true value. How far is the price below true value?
\end{exercise}

\begin{exercise}[$\star$]\label{exo:s2:systematic-credit-and-bond-etf-arbitrage:2}
The price is 300 basis points below true value and selling the basket costs 150. What does redeeming earn?
\end{exercise}

\begin{exercise}[$\star$]\label{exo:s2:systematic-credit-and-bond-etf-arbitrage:3}
Market spreads widen 150 basis points on bonds of average spread duration 6.5. Roughly how much does the basket lose?
\end{exercise}

\begin{exercise}[$\star\star$]\label{exo:s2:systematic-credit-and-bond-etf-arbitrage:4}
Why must \omterm{def:s2:systematic-credit-and-bond-etf-arbitrage:factor}{credit factors} be compared across similar durations?
\end{exercise}

\begin{exercise}[$\star\star$]\label{exo:s2:systematic-credit-and-bond-etf-arbitrage:5}
Why were \omterm{def:s2:systematic-credit-and-bond-etf-arbitrage:discount}{ETF discounts} larger for safer bonds in March 2020?
\end{exercise}

\begin{exercise}[$\star\star$]\label{exo:s2:systematic-credit-and-bond-etf-arbitrage:6}
Why does a bond ETF choose a redemption basket that is not the index?
\end{exercise}

\begin{exercise}[$\star\star\star$]\label{exo:s2:systematic-credit-and-bond-etf-arbitrage:7}
\emph{Coding.} Rerun with \texttt{SysCreditConfig(trade\_prob=0.05)}. What happens to the stale part of the discount at the stress's worst?
\end{exercise}

\begin{exercise}[$\star\star\star$]\label{exo:s2:systematic-credit-and-bond-etf-arbitrage:8}
\emph{Find the flaw.} ``The ETF trades 5\% below NAV, so buying it earns 5\% when the discount closes.''
\end{exercise}

\section{Problem: Which Price Is Right}

\begin{problem}[{Weekend problem --- credit factors and bond ETFs}]\label{pb:s2:systematic-credit-and-bond-etf-arbitrage:1}
The chapter's synthetic bonds and ETF and the public record.

\medskip\noindent\textbf{Part I --- Factors.}
\begin{enumerate}
\item Define a \omterm{def:s2:systematic-credit-and-bond-etf-arbitrage:factor}{credit factor}.
\item What did Houweling and van Zundert find?
\item Describe the synthetic bonds and the three signals.
\item Give the factors' Sharpe ratios and correlations.
\end{enumerate}

\medskip\noindent\textbf{Part II --- ETFs.}
\begin{enumerate}[resume]
\item Define create-redeem arbitrage and the \omterm{def:s2:systematic-credit-and-bond-etf-arbitrage:discount}{ETF discount}.
\item What did Koont, Ma, Pastor and Zeng find about baskets?
\item How is the synthetic NAV computed?
\item What is a portfolio trade?
\end{enumerate}

\medskip\noindent\textbf{Part III --- The stress.}
\begin{enumerate}[resume]
\item Describe the planted stress.
\item Decompose the worst discount.
\item When could the authorised participant profit?
\item What did Haddad, Moreira and Muir find?
\end{enumerate}

\medskip\noindent\textbf{Part IV --- The verdict.}
\begin{enumerate}[resume]
\item State the \emph{named result}: the \omterm{def:s2:systematic-credit-and-bond-etf-arbitrage:factor}{credit factors}' Sharpe ratios and the ETF arbitrage's return in the planted stress.
\item Which price is right, and why can no one see it?
\item Why is momentum weak here?
\item How would you estimate true value in a stress?
\item How would you backtest the ETF arbitrage honestly?
\item Which strategy file needs authorised-participant status?
\item How does this chapter relate to chapter~18?
\item In one sentence: what does a bond \omterm{def:s2:systematic-credit-and-bond-etf-arbitrage:discount}{ETF discount} measure?
\end{enumerate}
\end{problem}

\section{Interview questions}

\begin{interviewq}[$\star$ \iqroles{trader}]\label{iq:s2:systematic-credit-and-bond-etf-arbitrage:1}
How does creation and redemption keep an ETF near its NAV?
\end{interviewq}

\begin{interviewq}[$\star\star$ \iqroles{researcher}]\label{iq:s2:systematic-credit-and-bond-etf-arbitrage:2}
How would you build a value factor for corporate bonds?
\end{interviewq}

\begin{interviewq}[$\star\star$ \iqroles{trader}]\label{iq:s2:systematic-credit-and-bond-etf-arbitrage:3}
A client wants to sell a list of 200 bonds at once. How do you price it?
\end{interviewq}

\begin{interviewq}[$\star\star$ \iqroles{risk}]\label{iq:s2:systematic-credit-and-bond-etf-arbitrage:4}
Your book holds bond ETFs against their bonds. What happens in a stress like March 2020?
\end{interviewq}

\begin{interviewq}[$\star\star$ \iqroles{developer}]\label{iq:s2:systematic-credit-and-bond-etf-arbitrage:5}
Design a fair-value estimate for bonds that have not traded today.
\end{interviewq}

\begin{interviewq}[$\star\star\star$ \iqroles{researcher}]\label{iq:s2:systematic-credit-and-bond-etf-arbitrage:6}
If each bond trades on a share $p$ of days and true values follow a random walk, show that the NAV lags true value by $(1 - p)/p$ days on average,
and find the stale gap after a steady move of $m$ a day.
\end{interviewq}
